\begin{thebibliography}{99}
\small\raggedright
\bibitem{Parker2009} Parker J.S., Mullins M., Cheang M.C.U., et al. (2009). Supervised risk predictor of breast cancer based on intrinsic subtypes. \textit{Journal of Clinical Oncology} 27:1160--1167. \url{https://doi.org/10.1200/JCO.2008.18.1370}.

\bibitem{Mohammed2015} Mohammed H., Russell I.A., Stark R., et al. (2015). Progesterone receptor modulates estrogen receptor-alpha action in breast cancer. \textit{Nature} 523:313--317. \url{https://doi.org/10.1038/nature14583}.

\bibitem{Beleut2010} Beleut M., Rajaram R.D., Caikovski M., et al. (2010). Two distinct mechanisms underlie progesterone-induced proliferation in the mammary gland. \textit{Proceedings of the National Academy of Sciences} 107:2989--2994. \url{https://doi.org/10.1073/pnas.0915148107}.

\bibitem{Schramek2010} Schramek D., Leibbrandt A., Sigl V., et al. (2010). Osteoclast differentiation factor RANKL controls development of progestin-driven mammary cancer. \textit{Nature} 468:98--102. \url{https://doi.org/10.1038/nature09387}.

\bibitem{Toy2013} Toy W., Shen Y., Won H., et al. (2013). ESR1 ligand binding domain mutations in hormone-resistant breast cancer. \textit{Nature Genetics} 45:1439--1445. \url{https://doi.org/10.1038/ng.2822}.

\bibitem{Robinson2013} Robinson D.R., Wu Y., Vats P., et al. (2013). Activating ESR1 mutations in hormone-resistant metastatic breast cancer. \textit{Nature Genetics} 45:1446--1451. \url{https://doi.org/10.1038/ng.2823}.

\bibitem{Gruosso2016} Gruosso T., Mieulet V., Cardon M., et al. (2016). Chronic oxidative stress promotes H2AX protein degradation and enhances chemosensitivity in breast cancer patients. \textit{EMBO Molecular Medicine} 8:527--549. \url{https://doi.org/10.15252/emmm.201505891}.

\bibitem{GEO} NCBI Gene Expression Omnibus (2016). GSE45827: Expression data from breast cancer subtypes. \url{https://www.ncbi.nlm.nih.gov/geo/query/acc.cgi?acc=GSE45827}. Accessed 16 October 2026.

\bibitem{Smyth2004} Smyth G.K. (2004). Linear models and empirical Bayes methods for assessing differential expression in microarray experiments. \textit{Statistical Applications in Genetics and Molecular Biology} 3. \url{https://doi.org/10.2202/1544-6115.1027}.

\bibitem{Ritchie2015} Ritchie M.E., Phipson B., Wu D., et al. (2015). limma powers differential expression analyses for RNA-sequencing and microarray studies. \textit{Nucleic Acids Research} 43:e47. \url{https://doi.org/10.1093/nar/gkv007}.

\bibitem{BH1995} Benjamini Y., Hochberg Y. (1995). Controlling the false discovery rate: a practical and powerful approach to multiple testing. \textit{Journal of the Royal Statistical Society: Series B (Methodological)} 57:289--300. \url{https://doi.org/10.1111/j.2517-6161.1995.tb02031.x}.

\bibitem{Reactome2024} Milacic M., Beavers D., Conley P., et al. (2024). The Reactome Pathway Knowledgebase 2024. \textit{Nucleic Acids Research} 52:D672--D678. \url{https://doi.org/10.1093/nar/gkad1025}.

\bibitem{NCIHormone} National Cancer Institute (2026). Hormone therapy for breast cancer. \url{https://www.cancer.gov/types/breast/treatment/hormone-therapy}. Accessed 16 October 2026.

\bibitem{NCITargeted} National Cancer Institute (2026). Targeted therapy for breast cancer. \url{https://www.cancer.gov/types/breast/treatment/targeted-therapy}. Accessed 16 October 2026.

\bibitem{NCIDrugs} National Cancer Institute (2026). Drugs approved for breast cancer. \url{https://www.cancer.gov/about-cancer/treatment/drugs/breast}. Accessed 16 October 2026.

\bibitem{Cottu2018} Cottu P.H., Bonneterre J., Varga A., et al. (2018). Phase I study of onapristone, a type I antiprogestin, in female patients with previously treated recurrent or metastatic progesterone receptor-expressing cancers. \textit{PLOS ONE}. \url{https://doi.org/10.1371/journal.pone.0204973}.
\end{thebibliography}
